% =====================================================================
%  Formulación del outlet coronario no-Newtoniano de dos compartimentos
%  (updateRCRNonNew + outletFlow + calibración de Murray)
%  Archivo: CoupledLV0DCoronary3D.cpp
%  Copiar y pegar. Requiere amsmath, amssymb.
% =====================================================================
\documentclass[11pt]{article}
\usepackage{amsmath,amssymb}
\usepackage[margin=2.4cm]{geometry}
\newcommand{\dd}[2]{\frac{\partial #1}{\partial #2}}
\begin{document}

\section*{1. Ley de outlet no-Newtoniana (tubo de Carreau--Yasuda)}

Sangre de Carreau--Yasuda, viscosidad dependiente del cizallamiento $\dot\gamma$:
\begin{equation}
  \mu(\dot\gamma)=\mu_\infty+(\mu_0-\mu_\infty)
  \Bigl[1+(\lambda\dot\gamma)^{a}\Bigr]^{\frac{n-1}{a}} .
  \label{eq:cy}
\end{equation}
La \textbf{patología} se codifica cambiando el conjunto completo
$(\mu_0,\mu_\infty,\lambda,n,a)$, no solo $\mu_0$: al mover $\lambda,n,a$ se
desplaza toda la curva $\mu(\dot\gamma)$ y por tanto la resistencia efectiva en
\emph{todo} el rango de cizallamiento.

Para un tubo recto de radio $r$ y longitud $L$ con caída $\Delta p$, el esfuerzo
de corte en pared y el cizallamiento en pared son
\begin{equation}
  \tau_w=\frac{r\,|\Delta p|}{2L},
  \qquad
  \tau(\dot\gamma_w)=\mu(\dot\gamma_w)\,\dot\gamma_w=\tau_w
  \;\;\Rightarrow\;\;\dot\gamma_w .
  \label{eq:wall}
\end{equation}
El caudal se obtiene por la relación de Weissenberg--Rabinowitsch--Mooney
(cambiando la variable $\tau=\dot\gamma\,\mu(\dot\gamma)$,
$d\tau=(\mu+\dot\gamma\mu')\,d\dot\gamma$):
\begin{equation}
  Q(\Delta p)=\frac{\pi r^3}{\tau_w^{3}}
  \int_{0}^{\tau_w}\!\tau^2\,\dot\gamma(\tau)\,d\tau
  =\frac{\pi r^3}{\tau_w^{3}}
  \int_{0}^{\dot\gamma_w}\!\dot\gamma^{3}\,\mu^2\,
  \bigl(\mu+\dot\gamma\,\mu'\bigr)\,d\dot\gamma ,
  \label{eq:wrms}
\end{equation}
con derivada (usada en los solvers, \emph{exacta}):
\begin{equation}
  \boxed{\;\dd{Q}{\Delta p}=\frac{\pi r^3\,\dot\gamma_w-3\,Q}{\Delta p}\;}
  \qquad\text{(se sigue de }\tau_w=\tfrac{r}{2L}\Delta p\text{ en }\eqref{eq:wrms}).
  \label{eq:dQ}
\end{equation}
Notación: $q=f(\Delta p;L,r)$ con $q'=\partial f/\partial\Delta p$ de
\eqref{eq:dQ}. Fallback de Poiseuille $q=\Delta p/R_0$,
$R_0=8\mu_0 L/(\pi r^4)$ cuando $|\Delta p|$ es despreciable.

\section*{2. Modelo de outlet de dos compartimentos (por rama)}

Nodo arterial (subendocárdico) $p_c$ con compliancia $C$, comprimido por la
presión intramiocárdica $p_{im}$; nodo venoso $p_{ven}$ con compliancia $C_v$,
comprimido por $p_{im,v}$. El capacitor almacena la carga \emph{transmural}
$C(p_c-p_{im})$, de donde nace la bomba intramiocárdica.

\paragraph{Presiones tisulares (fuente cavitaria CIEP + activa SIP), filtradas.}
\begin{align}
  p_{im}^{\text{tgt}}   &= \alpha\,p_{LV}+\alpha_a\,\tau_c, &
  p_{im,v}^{\text{tgt}} &= \alpha_v\,p_{LV}+\alpha_{v,a}\,\tau_c,
  \label{eq:tgt}\\
  p_{im}^{n} &= p_{im}^{n-1}+\theta\bigl(p_{im}^{\text{tgt}}-p_{im}^{n-1}\bigr),
  & \theta &= \frac{\Delta t}{\Delta t+\tau_f},
  \label{eq:filt}
\end{align}
($\tau_c$ = estrés activo de fibra del modelo 0D). El filtro
\eqref{eq:filt} acota $\dot p_{im}$ y elimina la aguja de un paso de las fases
isovolumétricas sin borrar la bomba. Las fuentes de bomba son
\begin{equation}
  Q_{im}=C\,\frac{p_{im}^{n}-p_{im}^{n-1}}{\Delta t},
  \qquad
  Q_{im,v}=C_v\,\frac{p_{im,v}^{n}-p_{im,v}^{n-1}}{\Delta t}.
  \label{eq:pump}
\end{equation}

\paragraph{Ramales.} Proximal (epicárdico) y distal no-Newtonianos por
\eqref{eq:wrms}; drenaje venoso al atrio derecho $P_{RA}$ \textbf{lineal}:
\begin{equation}
  q_d=f\!\bigl(p_c-p_{ven};L_D,r_D\bigr),\qquad
  q_v=\frac{p_{ven}-P_{RA}}{R_v},\qquad
  p_{out}=p_c+\Delta p_P(Q).
  \label{eq:branches}
\end{equation}
($\Delta p_P(Q)$ = inversión de \eqref{eq:wrms} en el ramal proximal para el
caudal $Q$ que entrega la malla 3D.)

\paragraph{Balance de nodo (almacenamiento transmural).}
\begin{equation}
  C\,\frac{d}{dt}(p_c-p_{im})=Q-q_d,
  \qquad
  C_v\,\frac{d}{dt}(p_{ven}-p_{im,v})=q_d-q_v .
  \label{eq:balance}
\end{equation}

\section*{3. Discretización implícita y Newton $2\times2$}

Con Euler implícito, incógnitas $(p_c,p_{ven})$:
\begin{align}
  R_1 &= \frac{C}{\Delta t}\bigl(p_c-p_c^{\,n}\bigr)-Q_{im}+q_d-Q=0,
  \label{eq:R1}\\
  R_2 &= \frac{C_v}{\Delta t}\bigl(p_{ven}-p_{ven}^{\,n}\bigr)-Q_{im,v}-q_d+q_v=0.
  \label{eq:R2}
\end{align}
Con $q_d'=\partial q_d/\partial(p_c-p_{ven})$ y $q_v'=1/R_v$, el Jacobiano es
\begin{equation}
  J=\begin{pmatrix}
    \dfrac{C}{\Delta t}+q_d' & -q_d' \\[6pt]
    -q_d' & \dfrac{C_v}{\Delta t}+q_d'+q_v'
  \end{pmatrix}.
  \label{eq:jac}
\end{equation}

\section*{4. Calibración newtoniana con geometría fija (ley de Murray)}

\emph{Se calibra una sola vez con una viscosidad newtoniana $\mu_N$ para fijar
la geometría de forma imparcial}; luego el paso temporal usa la reología real
\eqref{eq:cy}. Así, entre patologías cambia solo $\mu(\dot\gamma)$, nunca la
geometría. Áreas de salida $A_i$ medidas de la malla, radio equivalente
$r_i=\sqrt{A_i/\pi}$, reparto por Murray ($Q\propto r^3$):
\begin{align}
  w_i &= \frac{r_i^{3}}{\sum_j r_j^{3}}, &
  Q_i &= Q_{\text{tot}}\,w_i, &
  \Delta P &= p_{ar}(0)-\alpha\,p_{LV}(0),\\
  R_i &= \frac{\Delta P}{Q_i}, &
  R_d &= (1-f_p)R_i, & R_p &= f_p R_i,
  \label{eq:murray}
\end{align}
radios efectivos por inversión de Poiseuille (con $\mu_N$) y compliancias:
\begin{equation}
  r_D=\Bigl(\frac{8\mu_N L_D}{\pi R_d}\Bigr)^{1/4},\quad
  r_P=\Bigl(\frac{8\mu_N L_P}{\pi R_p}\Bigr)^{1/4},\quad
  C=\frac{\tau_{RCR}}{R_d},\quad
  C_v=\kappa_v C,\quad R_v=\rho_v R_d .
  \label{eq:geom}
\end{equation}

\section*{5. Estado de la formulación}

\begin{itemize}
\item \textbf{Ley de tubo \eqref{eq:wrms}--\eqref{eq:dQ}: correcta.} Es
  Rabinowitsch/WRMS con Jacobiano analítico exacto; verificado por
  diferenciación directa.
\item \textbf{Calibración newtoniana \eqref{eq:murray}--\eqref{eq:geom}:
  correcta y recomendada.} Geometría fija $\Rightarrow$ comparación imparcial
  entre reologías. La sensibilidad a la patología entra por $(\mu_0,\mu_\infty,
  \lambda,n,a)$ en \eqref{eq:cy}.
\item \textbf{Bomba \eqref{eq:tgt}--\eqref{eq:pump} y balance
  \eqref{eq:balance}: correctos.} Reproducen impedimento sistólico del flujo
  arterial y aumento sistólico del retorno venoso.
\item \textbf{Correcciones aplicadas:}
  (i) drenaje venoso $q_v=(p_{ven}-P_{RA})/R_v$ \emph{lineal} (antes se pasaba
  la resistencia $R_v$ como radio a la ley de tubo, clavando $p_{ven}\approx
  P_{RA}$ y anulando el compartimento venoso);
  (ii) $\Delta P$ de calibración con $\alpha=$ \texttt{intramyocardialFraction}
  (antes $0.45$ hardcodeado, inconsistente con $0.65$ del modelo);
  (iii) etiquetas de salida del CSV a los atributos reales $\{7,8,9,10,14,15\}$.
\item \textbf{Perillas físicas (no errores):} la retirada protodiastólica la
  fija $\tau_{RCR}=R_dC$ (arterial) y $\tau_v=R_vC_v=\rho_v\kappa_v\tau_{RCR}$
  (venosa); subirlas la hace más lenta.
\end{itemize}

\end{document}
